\begin{tikzpicture}[
    node distance=16mm and 18mm,
    >=Stealth,
    mainblock/.style={
        draw=black!85, 
        line width=1pt, 
        rounded corners=3pt, 
        minimum width=38mm, 
        minimum height=14mm, 
        align=center, 
        fill=gray!8, 
        inner sep=4pt
    },
    opsblock/.style={
        mainblock,
        fill=white,
        draw=black!70,
        line width=0.75pt,
        dashed
    },
    flow/.style={->, line width=1pt, draw=black!90},
    feedback/.style={->, line width=0.75pt, dashed, draw=black!70},
    boundary/.style={
        draw=black!40, 
        line width=0.5pt, 
        dashed, 
        rounded corners=4pt, 
        fill=gray!2, 
        inner sep=18pt
    },
    edgelabel/.style={
        font=\scriptsize\sffamily, 
        text=black!75, 
        fill=gray!2, 
        inner sep=2pt
    }
]

\node[mainblock] (sched) {\textbf{调度与流量整形}\\[1.5mm]\footnotesize 令牌桶 / SLO};
\node[mainblock, left=of sched] (request) {\textbf{请求进入}\\[1.5mm]\footnotesize 服务接口};
\node[mainblock, right=of sched] (exec) {\textbf{执行与返回}\\[1.5mm]\footnotesize 流式输出};

\node[opsblock, below=18mm of sched] (ops) {\textbf{观测与运维}\\[1.5mm]\footnotesize 日志采集 $\vert$ 性能监控 $\vert$ 异常回退};

\path (exec.east) ++(28mm, 0) coordinate (right_padding);

\begin{scope}[on background layer]
    \node[boundary, fit=(request) (exec) (ops) (sched) (right_padding)] (system_boundary) {};
\end{scope}

\node[draw=black!40, line width=0.5pt, fill=white, rounded corners=2pt, inner sep=4pt, 
      anchor=south east, xshift=-4mm, yshift=4mm] at (system_boundary.south east) {
    \begin{tabular}{@{}cl@{}}
    \tikz[baseline=-0.5ex]\draw[flow] (0,0) -- (0.4,0); & \scriptsize 主请求流 \\
    \tikz[baseline=-0.5ex]\draw[feedback] (0,0) -- (0.4,0); & \scriptsize 运维反馈流
    \end{tabular}
};

\draw[flow] (request) -- (sched);
\draw[flow] (sched) -- (exec);
\draw[feedback] (exec.south) |- node[pos=0.75, edgelabel, anchor=south] {监控数据采集} (ops.east);
\draw[feedback] (ops.west) -| node[pos=0.25, edgelabel, anchor=south] {异常回退 / 配置调整} (request.south);

\end{tikzpicture}